% Burnett's "stress cycle" with his new-boss example. Steps reveal on slides 1-5, the loop closes on slide 6.
\begin{tikzpicture}[st/.style={flow=3.15cm, font=\small, inner sep=3pt}]
  \node[st, fill=butter!55!cardwhite] (n1) at (0,1.75) {\emoji{boss}\ Նոր պետ, անհնար ծավալ};
  \node[st, fill=peach!45!cardwhite, visible on=<2->] (n2) at (3.15,0.5) {\emoji{hourglass}\ Երկար ժամեր, քրոնիկ սթրես};
  \node[st, fill=peach!65!cardwhite, visible on=<3->] (n3) at (1.95,-1.45) {\emoji{burger}\emoji{beer}\ Արագ սնունդ ու ալկոհոլ՝ «հանգստանալու»};
  \node[st, fill=brainpink!40!cardwhite, visible on=<4->] (n4) at (-1.95,-1.45) {\emoji{weary}\ Առողջությունն ու տրամադրությունը վատանում են};
  \node[st, fill=brainpink!60!cardwhite, visible on=<5->] (n5) at (-3.15,0.5) {\emoji{batterylow}\ Ավելի խոցելի ես նոր սթրեսի համար};
  \draw[arr, visible on=<2->] (n1.east) to[bend left=25] (n2.north);
  \draw[arr, visible on=<3->] (n2.south) to[bend left=20] (n3.east);
  \draw[arr, visible on=<4->] (n3.west) to[bend left=15] (n4.east);
  \draw[arr, visible on=<5->] (n4.west) to[bend left=20] (n5.south);
  \draw[arr=pinkdark, visible on=<6->] (n5.north) to[bend left=25] (n1.west);
  \node[font=\blackfont\small, text=pinkdark, visible on=<6->] at (0,0.1) {\emoji{cyclone}\ ու նորից};
\end{tikzpicture}
